\newpage
\section{Background Knowledge and Technologies}

This chapter assembles the evidence on which the proposed design rests. Sections 2.1 and 2.2 concern learning and measurement: whether performing a task conveys understanding that description cannot, and whether competence can be inferred from performance. Section 2.3 examines game-based delivery, including evidence qualifying its effectiveness. Section 2.4 summarises career-development theory. Sections 2.5 and 2.6 address the generative and assessment technologies and the data resources available for grounding and evaluation, and Section 2.7 considers responsibility constraints arising from the intended user population.

\subsection{Experiential learning and realistic job previews}

Kolb's model of experiential learning\cite{r12} characterises learning as a cycle in which concrete experience is transformed through reflective observation into abstract conceptualisation and subsequent active experimentation. The model's implication for occupational information is direct: an occupational description supplies conceptual content without the experience from which, on this account, durable understanding is constructed. Whatever a description conveys, it does not engage the transformation process the model describes.

The realistic job preview literature provides the corresponding empirical tradition\cite{r13}. A realistic job preview presents a balanced account of a role, including features an applicant might find unattractive, in place of a promotional account. Four meta-analyses converge on its effects\cite{r8,r9,r14,r15}. Phillips\cite{r8}, synthesising forty studies, associated such previews with reduced inflation of expectations, increased self-selection, improved performance and lower voluntary turnover, and additionally found richer audiovisual previews more effective than written ones --- a finding that bears directly on delivery format. Earnest, Allen and Landis\cite{r9}, in a path analysis over fifty-two studies and approximately seventeen thousand participants, identified perceived honesty as the mediating mechanism, indicating that the effect depends on the preview's candour rather than merely on its informational content. No meta-analysis has since superseded it; the most recent umbrella review of the field identifies it as the current synthesis\cite{r16}.

Whether simulated experience transfers to real performance is addressed by the simulation-based learning literature. Cook et al.\cite{r17} and McGaghie et al.\cite{r18} established substantial effects for technology-enhanced simulation within health professions education, where the paradigm originated. Of greater relevance here, Chernikova et al.\cite{r19}, synthesising 145 studies in higher education, reported a large overall effect that generalised across academic domains rather than remaining confined to clinical training, and further found that scaffolding --- worked examples for novices, reflective prompts for more advanced learners --- moderated outcomes.

Taken together, these three literatures support a specific proposition: that an honest, performed representation of a role conveys occupational information more effectively than description, and that the effect is not restricted to the clinical domains in which simulation research began.

\subsection{Assessment of competence from behaviour, and competency-based progression}

If a learner performs tasks rather than answering questions about themselves, the assessment problem changes accordingly. Evidence-centered design, developed for assessment within game environments under the term stealth assessment\cite{r20}, addresses this by decomposing the problem into three linked models: a competency model specifying what is claimed to be measured, an evidence model specifying which observable behaviours constitute evidence for those claims and with what weight, and a task model specifying situations that elicit the relevant behaviour.

{\small
\begin{longtable}{|>{\raggedright\arraybackslash}p{\dimexpr 0.159\textwidth-2\tabcolsep-\arrayrulewidth\relax}|>{\raggedright\arraybackslash}p{\dimexpr 0.317\textwidth-2\tabcolsep-\arrayrulewidth\relax}|>{\raggedright\arraybackslash}p{\dimexpr 0.519\textwidth-2\tabcolsep-\arrayrulewidth\relax}|}
\caption{Evidence-centered design models and the requirements they impose.} \\
\hline
\textbf{Model} & \textbf{Function} & \textbf{Requirement it imposes} \\ \hline
\endfirsthead
\hline
\textbf{Model} & \textbf{Function} & \textbf{Requirement it imposes} \\ \hline
\endhead
Competency & Defines the constructs claimed to be measured & Constructs must be specified independently of any single task if inferences are to transfer \\ \hline
Evidence & Maps observed behaviour to inferences about the constructs & Each observable must be attributable to a named construct; unattributable observations yield uninterpretable scores \\ \hline
Task & Specifies situations that elicit relevant behaviour & The target construct must be necessary to task completion, not merely permitted by it \\ \hline
\end{longtable}
}

The third requirement is the most demanding and the most frequently unmet. Shute et al.\cite{r21} obtained significant correlations between in-game estimates of problem-solving and external measures precisely because the tasks could not be completed without deploying the measured competency. Where a task admits a route that bypasses the target construct, the evidence model measures something other than what it claims. Task authoring in such systems is consequently a psychometric activity rather than a purely editorial one.

A related question is what an assessment result entitles a learner to do next. Competency-based progression\cite{r22} treats competence as a transferable and independently verifiable unit rather than as an attribute of a completed course or an occupied role. On this account, competence demonstrated in one context constitutes evidence bearing on readiness for adjacent contexts. Establishing which contexts are genuinely adjacent requires an external basis, and standard occupational taxonomies supply one by describing occupations in terms of constituent skills and work activities, discussed in Section 2.6.

Where a task asks the learner to choose among actions rather than to write, the relevant assessment tradition is the situational judgment test, which presents work situations with response options that differ in effectiveness and is commonly characterised as a low-fidelity simulation\cite{r23}. A meta-analysis of 102 validity coefficients from 10,640 people found situational judgment tests validly and generalisably related to job performance ($\rho$ = .34)\cite{r24}, and their items measure distinguishable constructs such as interpersonal skill, teamwork, and leadership rather than a single ability\cite{r25}. Two further findings bear on design. The established development process derives scenarios and response options from critical incidents supplied by subject-matter experts\cite{r26}, and a separate group of experts rates each option's effectiveness to form the scoring key\cite{r27}. The wording of the response instruction also changes the construct measured: asking what a respondent would do relates more to personality, whereas asking which action is most effective relates more to cognitive ability\cite{r28}. Options scored in this way are keyed along a range from ineffective or harmful to highly effective, and each carries an expert-rated profile of effectiveness across competencies. An option may accordingly be correct, in the sense of the most effective response; a trade-off, effective for one competency at a cost to another; or wrong, ineffective or harmful and keyed with negative effectiveness.

\subsection{Game-based delivery: evidence, moderators, and limits}

Game-based delivery is frequently proposed for educational systems and frequently overstated. The evidence supports it conditionally, and the conditions are specific enough to constrain design.

Effects and moderators. Hamari, Koivisto and Sarsa\cite{r29} found gamification effects generally positive but dependent on context and user. Sailer and Homner\cite{r30}, in the principal learning-specific synthesis, reported significant but modest effects --- cognitive g = .49, motivational g = .36, behavioural g = .25 --- and identified narrative and the combination of competition with collaboration as moderators, noting the effectiveness of narrative paired with a persistent developing avatar. Bai, Hew and Huang\cite{r31} and Clark, Tanner-Smith and Killingsworth\cite{r32} report comparable medium effects with design features moderating outcomes.

The competence qualification. A meta-analysis by Li, Hew and Du\cite{r33} disaggregates these effects by psychological need and yields a result that materially constrains claims made for such systems: gamification produced large gains in perceived autonomy (g = 0.638) and relatedness (g = 1.776) but only a marginal effect on competence (g = 0.277). Gamification, in other words, reliably improves how learners feel about an activity while doing comparatively little for what they can do. Zeng et al.\cite{r34} add that particular element combinations underperform, with one common configuration associating negatively with academic performance. Any system asserting skill development must therefore locate the mechanism of that development in something other than reward structures.

Evidence by genre. Genre-level findings differentiate more usefully than aggregate ones. Role-play methods show an overall effect of approximately 0.82 with the largest component on skill acquisition\cite{r35}, the strongest genre-level support for role-based designs. Simulation games carry the most occupationally relevant evidence and the clearest caveat: Sitzmann\cite{r36}, synthesising sixty-five studies, found simulation-taught trainees exceeded comparison groups on self-efficacy, declarative and procedural knowledge and retention, but only where the game actively conveyed material, permitted learner control, and supplemented other instruction --- when it constituted the sole method, comparison groups performed better, and the analysis reported evidence of publication bias. Wouters et al.\cite{r37} similarly found serious games effective for learning and retention while not necessarily more motivating than conventional instruction, and most effective as a supplement across multiple sessions. Narrative and branching-story formats reinforce engagement and learning\cite{r38}, consistent with the narrative moderator above. Procedurally varied and randomised formats support replayability, but the evidence here is double-edged: randomised-reward mechanisms are robustly associated with problem gambling and reduced wellbeing among purchasers\cite{r39}, which distinguishes variation used for content variety from variable-ratio reward scheduling.

Motivational basis. These findings are commonly interpreted through self-determination theory\cite{r40}, under which motivation depends on satisfaction of needs for competence, autonomy and relatedness. The mapping is reasonably direct: progression and mastery structures address competence, meaningful choice addresses autonomy, and social or comparative features address relatedness, the least studied of the three in game research.

The implication for design is that engagement mechanisms and learning mechanisms should not be conflated. Narrative, progression and choice are supported as engagement structures; skill acquisition, on this evidence, must derive from the authenticity of the tasks themselves.

\subsection{Career-development theory}

Four theoretical traditions inform career guidance provision and the design of systems supporting it. Holland's typology\cite{r41} models occupational choice as person--environment congruence and underlies the interest inventories in widespread use, including the public-domain instrument distributed with the O*NET occupational database. Social cognitive career theory\cite{r10,r42} identifies self-efficacy, built principally through mastery experience, as a determinant of interest formation and choice --- a mechanism accessible to systems in which users perform tasks rather than report preferences, and inaccessible to instruments that only elicit self-description. Super's life-span theory\cite{r43} holds that vocational self-concept develops over time, so that occupational fit is properly understood as changing rather than fixed. Krumboltz's happenstance learning theory\cite{r44} extends this by treating unplanned events as a normal and potentially productive component of career development rather than as disruption.

Social cognitive career theory also specifies how self-efficacy forms, which determines what a system of this kind can and cannot affect. Following Bandura\cite{r45}, self-efficacy derives from four sources: mastery experiences, vicarious experiences of others' performance, verbal and social persuasion, and affective and physiological states. Their contributions are unequal. A model-based meta-analysis identified mastery experience as the strongest source\cite{r11}; the finding has been extended to career exploration and decision-making\cite{r46} and confirmed longitudinally in a three-wave study of 1,512 university students, in which mastery experience exerted the greatest sustained influence on career decision self-efficacy\cite{r47}. Two qualifications follow for systems that deliver experience through simulated tasks. Such systems engage mastery experience directly and persuasion through feedback, but provide little vicarious experience and do not target affective states. And because self-efficacy is task-specific, a system can build it only for the tasks it actually presents: where interaction is text-based, it can address competencies expressed through written analysis, prioritisation, and decision-making, but not those expressed through speech or physical performance, such as oral fluency in an interview. Section 4.2.3 compares these competencies with their counterparts in live work situations.

The last two traditions jointly support a position that career guidance systems seldom represent: that a career decision is situated within a whole life, subject to circumstances --- obligations, health, relationships, relocation, changing values --- that alter what constitutes a suitable path. Where a system models occupational fit as a fixed verdict, this dimension is lost. The consideration carries particular weight in contexts where career decisions are made within strong family and social expectations, as the Vietnamese survey data indicate\cite{r4}, and it provides the theoretical basis for treating work--life balance as an attribute of an occupation rather than as an externality.

The limitations of self-report instruments follow from the same body of theory. Beyond the contested stability of type classifications noted in Section 1.1\cite{r6}, self-report is subject to social desirability effects, and congruence measures derived from interest typologies correlate only modestly with subsequent satisfaction. These are not implementation defects but properties of the method: asking what a person believes they would enjoy produces evidence of a different kind from observing what they do.

Frameworks of career decision difficulty. The taxonomy of Gati, Krausz and Osipow\cite{r1}, on which Section 1.1 relies, is one of several frameworks describing why career decisions stall, and it is among the oldest still in use, so its choice requires justification. Table 2.2 compares it with five others.

{\small
\begin{longtable}{|>{\raggedright\arraybackslash}p{\dimexpr 0.193\textwidth-2\tabcolsep-\arrayrulewidth\relax}|>{\raggedright\arraybackslash}p{\dimexpr 0.092\textwidth-2\tabcolsep-\arrayrulewidth\relax}|>{\raggedright\arraybackslash}p{\dimexpr 0.193\textwidth-2\tabcolsep-\arrayrulewidth\relax}|>{\raggedright\arraybackslash}p{\dimexpr 0.307\textwidth-2\tabcolsep-\arrayrulewidth\relax}|>{\raggedright\arraybackslash}p{\dimexpr 0.209\textwidth-2\tabcolsep-\arrayrulewidth\relax}|}
\caption{Frameworks of career decision difficulty compared with the taxonomy of Gati, Krausz and Osipow.} \\
\hline
\textbf{Framework} & \textbf{Year} & \textbf{What it describes} & \textbf{Structure} & \textbf{Missing occupational information as a separate component} \\ \hline
\endfirsthead
\hline
\textbf{Framework} & \textbf{Year} & \textbf{What it describes} & \textbf{Structure} & \textbf{Missing occupational information as a separate component} \\ \hline
\endhead
Taxonomy of career decision-making difficulties (CDDQ)\cite{r1} & 1996 & Difficulties arising before and during the decision & Ten categories in three clusters: lack of readiness, lack of information, inconsistent information & Yes: within the lack-of-information cluster, alongside information about the self, the process, and ways of obtaining information \\ \hline
Emotional and personality-related career decision-making difficulties (EPCD)\cite{r106} & 2008 & Emotional and personality-related difficulties & Three clusters: pessimistic views, anxiety, self-concept and identity & No \\ \hline
Career Indecision Profile\cite{r107} & 2013 & Sources of indecision, including personality & Four factors: neuroticism and negative affectivity, choice and commitment anxiety, lack of readiness, interpersonal conflicts & No separate component \\ \hline
Career Decision Self-Efficacy scale\cite{r108} & 1983 & Confidence in carrying out decision tasks & Five subscales: self-appraisal, occupational information, goal selection, planning, problem solving & Partly: confidence in gathering information, not the information that is lacking \\ \hline
Career Thoughts Inventory\cite{r109} & 1996 & Dysfunctional career thoughts & Three scales: decision-making confusion, commitment anxiety, external conflict & No \\ \hline
Career Decision Scale\cite{r110} & 1976 & Certainty and indecision & A certainty scale and an indecision scale & No separate component \\ \hline
\end{longtable}
}

The comparison shows a division of labour rather than a succession. The later frameworks do not replace the taxonomy; they describe other sources of difficulty --- emotional and personality-related dispositions\cite{r106,r107,r109}, or confidence in carrying out decision tasks\cite{r108} --- and the first of them was developed with the taxonomy's own author to complement it\cite{r106}. Only the taxonomy isolates missing information about occupations as a separately measured category, distinct from readiness and from inconsistent information. That category is the difficulty a system supplying occupational information can address directly, whereas dispositions are unlikely to be changed by information alone. The taxonomy's age is offset by recent structural evidence: its structure was re-examined across thirteen countries in 2023\cite{r2}, a study that re-validates the taxonomy rather than proposing an alternative to it.

Applicability to Vietnam. The cross-national evidence concerns the taxonomy's structure. Levin et al.\cite{r2} compared factor models on thirty-nine samples (N = 19,562) in nine language versions, which the authors describe as covering thirteen countries, and found the original structure to fit best, including in the South Korean and Malaysian samples, the only Asian ones. Measurement invariance could be tested only within the English and the French versions; it held at the scalar level for the French version but not at the metric level for the English one, so scores are not established as comparable across countries. The lack-of-information cluster was the most reliable of the three and the lack-of-readiness cluster the weakest, which is also where earlier Taiwanese, Chinese, and Korean studies reported departures from the original structure\cite{r111,r112,r113}, one of them placing external conflicts, which include the expectations of family, with readiness\cite{r111}. No validated Vietnamese version of the questionnaire was identified. The project therefore relies on the taxonomy only for the category that replicates most consistently, missing information, and does not use the questionnaire's scores; using it as an outcome measure with Vietnamese users would first require local validation, with particular attention to readiness and to family-related external conflicts.

\subsection{Generative and assessment technologies}

Content generation. Automatic generation of educational content predates large language models but was long confined to shallow factual items\cite{r48}. Language models extended this substantially: Sarsa et al.\cite{r49} generated programming exercises with solutions, tests and explanations, a domain directly relevant to information technology scenarios, while noting that quality oversight remained necessary; Doughty et al.\cite{r50} generated items calibrated to specified cognitive levels. Controlled comparison indicates the outputs can be psychometrically adequate --- one study found model-authored assessment items statistically indistinguishable from human-authored items on difficulty, discrimination and reliability, with the effect most pronounced in computer science\cite{r51}. The consistent qualification across this literature concerns grounding: fluent generation is readily obtained, while factual fidelity to a specified domain requires explicit constraint, typically through retrieval or authored scaffolding\cite{r52}.

Interactive narrative and its control. Systems in which a language model conducts an interactive scenario draw on work demonstrating that such models can maintain memory, synthesise reflections and plan behaviour over extended interaction\cite{r53}, and can track game state while a human retains narrative authority\cite{r54,r55}. The same literature documents the principal failure mode, namely degradation of consistency over long horizons, and the mitigations that address it. Plan-driven generation frameworks report improved continuity and reduced spatiotemporal inconsistency relative to unconstrained prompting\cite{r56}, indicating that authored structure rather than model scale is the operative variable.

Automated assessment of open responses. The use of language models as evaluators has developed through three identifiable stages, each imposing a different requirement on systems that adopt it. Zheng et al.\cite{r57} established feasibility, reporting agreement with human preferences above eighty per cent --- comparable to inter-human agreement --- while documenting position, verbosity and self-preference biases in the same work. Liu et al.\cite{r58} addressed method, showing that structured chain-of-thought prompting against an explicit rubric improves alignment, while noting a residual preference for model-generated text. The third stage concerns validation and is the most consequential. A systematic evaluation of twenty-one evaluators across approximately 541,000 judgments\cite{r59} found that raw agreement systematically overstates chance-corrected performance, reporting an approximately forty-one point discrepancy between exact-match agreement and Cohen's kappa on one benchmark, and that high test--retest consistency can coexist with severe position bias. Dedicated benchmarks quantify the ceiling: on deliberately difficult response pairs with objective correctness labels, the strongest general-purpose evaluator without extended reasoning reached approximately sixty-four per cent accuracy, and models with extended reasoning seventy-five to eighty-one per cent\cite{r60}.

Reliability of rubric scoring by language models. Whether such evaluators are themselves reliable for rubric-based scoring of learners' constructed responses has been studied directly, and the evidence supports a conditional answer. A synthesis of sixty-five studies of essay scoring found agreement between language models and human raters to be highly context-dependent, with reported quadratic weighted kappa spanning nearly the whole range and many values below the 0.70 conventionally required\cite{r101}. Agreement approaches that between human raters under identifiable conditions: a strong model, an explicit rubric, and calibration examples. With few-shot prompting, GPT-4 marked short answers in school subjects at a Cohen's kappa of 0.70, against 0.75 between human markers\cite{r102}, and with calibration examples its ratings of short second-language essays approached those of dedicated scoring systems\cite{r103}. Without such support agreement is moderate: a small open model reached a mean quadratic weighted kappa of about 0.5 zero-shot on a standard essay corpus, against 0.79 for a fine-tuned scorer\cite{r104}. Consistency across repeated runs can be high, with intraclass correlations of 0.94 to 0.99 reported for GPT-4\cite{r105}, but consistency is not validity. Two implications follow. Judging models should be strong, anchored to an explicit rubric, and supplied with scored examples; and reliability cannot be inferred from the literature for a new task and domain, none of these studies concerning occupational tasks, but must be established against experts in each case.

Four requirements follow for any system adopting this technique. Agreement must be reported with chance correction rather than as exact match. The evaluator should be drawn from a different model family than the generator, since self-preference otherwise contaminates the loop. Option ordering and response length must be controlled, since both bias scores independently of quality. And because accuracy degrades unevenly with item difficulty, aggregate figures conceal failure on precisely the reasoning-intensive items such systems most need to assess, which requires stratified reporting.

Expert involvement in authoring. Maintaining accuracy in generated instructional content is a multi-party problem. Recent work converges on a layered pipeline: a curated source of truth, retrieval-grounded generation, expert review prior to publication, and a feedback loop from use. Evaluation of such pipelines indicates that automated checks reliably enforce verifiable properties while domain experts remain necessary for pedagogical adequacy and for the distinctions that make an item discriminating\cite{r61}. The division of labour this implies --- automated first-pass generation under expert validation --- is the current best-supported arrangement for content that must remain factually accurate about a real domain.

\subsection{Data resources and benchmarks}

Systems making claims about occupations and about assessment accuracy can be grounded in, and evaluated against, established resources rather than bespoke data collection.

Occupational and skills taxonomies. The O*NET database maintained by the United States Department of Labor describes approximately nine hundred occupations linked to more than nineteen thousand task statements across some two thousand detailed work activities\cite{r62}. The European ESCO classification describes 3,039 occupations and 13,939 skills across twenty-eight languages, with an official correspondence to O*NET\cite{r63}. Both express occupations as compositions of skills and activities, which supplies an external basis for claims about occupational adjacency and a source for authoring occupationally representative tasks.

Evaluation benchmarks. For automated assessment, MT-Bench and the associated arena methodology\cite{r57}, JudgeBench\cite{r60} and RewardBench\cite{r64} provide difficulty-calibrated reference sets and, more importantly, an established reporting methodology. For generated instructional content, EQGBench\cite{r65} and EduBench\cite{r66} provide protocols for evaluating generation quality. Text-based interactive environments including TextWorld\cite{r67} and Jericho\cite{r68} serve as established settings for evaluating agents in branching narrative.

Measurement practice. Two constructs require separate instruments and should not be reported as a single accuracy figure. Generation quality concerns whether produced content preserves required domain facts and elicits the intended constructs. Expert review of such content has an established quantitative form in content validation: experts rate each item's relevance on a four-point scale, and the item-level content validity index (I-CVI) is the proportion of experts rating it relevant. Three experts is the accepted minimum for a content-validation effort\cite{r69,r70}, and an I-CVI of .78 or higher with three or more experts, together with a scale-level average (S-CVI/Ave) of at least .90, indicates good content validity\cite{r70}. Assessment accuracy concerns agreement between automated and expert judgment of the same response. Chance-corrected agreement is required because raw agreement overstates performance, and because rubric scores are ordinal, quadratic weighted kappa is the appropriate statistic, penalising large disagreements more heavily than small ones. For automated scoring, Williamson, Xi and Breyer\cite{r71} set acceptance at a quadratic weighted kappa of at least 0.70 between automated and human scores, a standardised mean score difference of at most 0.15, and a degradation of no more than 0.10 from human--human to automated--human agreement. Agreement among the human raters themselves is assessed with Krippendorff's alpha, for which .800 indicates reliable ratings and .667 is the lower bound for tentative conclusions\cite{r72}. Agreement is reported stratified by item difficulty and alongside test--retest consistency. Two further techniques allow evaluation before human responses are available. Perturbation checklists alter a response along a single targeted criterion while holding the others constant, which supplies ground truth by construction and exposes whether an evaluator is sensitive to that criterion\cite{r73}; and a panel of judging models drawn from different model families has been found to outperform a single large judge, with less intra-model bias and at lower cost\cite{r74}, a result obtained for binary and pairwise judgments whose extension to ordinal rubric scoring remains to be tested. The distinction matters because the three failure modes --- biased judgment, unstable judgment, and judgment that fails selectively on difficult items --- are separately diagnosable only when reported separately. As to expected magnitude, the closest methodological precedent, a multi-agent architecture performing unobtrusive assessment in a serious game, reports correlations between behavioural estimates and external outcomes in the range of 0.28 to 0.33, with near-zero correlation against prior attainment\cite{r75}. The latter figure is as informative as the former, since it evidences discriminant validity. Moderate correlations are the realistic expectation for open-ended behavioural assessment.

\subsection{Responsibility considerations}

Three considerations arise from the intended user population and the nature of the data involved, each corresponding to a requirement of the project's specification.

Data protection for young users. A system addressed partly to secondary and tertiary students processes behavioural and performance data concerning individuals who may be minors, engaging Vietnam's Law on Personal Data Protection\cite{r99} and its implementing decree\cite{r100}, both in force from 1 January 2026, and comparable frameworks for users elsewhere. Under that law, processing data collected in Vietnam on a platform located abroad is a cross-border transfer of personal data, which bears on any system that sends users' responses to a model provider outside the country. Aggregated statistics derived from individual performance carry a residual privacy surface, since aggregate presentations can under some conditions permit inference about individuals.

Fairness in automated assessment. An automated judgment of a user's demonstrated competence may influence that user's estimation of their own capability. The evaluator biases documented in Section 2.5 are therefore not solely accuracy concerns but fairness concerns, and the corresponding controls --- chance-corrected validation, cross-family evaluation, rubric anchoring --- function as fairness safeguards. Presenting assessment formatively, and not as a determination affecting access to opportunity, keeps the consequences proportionate to the reliability the method can currently support.

Accuracy of occupational representation. The realistic job preview evidence carries an ethical as well as a pedagogical implication: a system's value depends on representing occupations candidly, including their disadvantages, and is undermined if representation is shaped by commercial interest. Expert validation of authored material, discussed in Section 2.5, is the principal safeguard against a system conveying an inaccurate picture of a real profession with apparent authority.

\subsection{Summary}

The background supports a coherent position. Performed, honest representation of occupational work conveys information that description does not, and this effect generalises across domains (2.1). Competence can be inferred from behaviour provided the assessment is constructed so that target constructs are necessary to task completion, and such competence is meaningfully transferable across related occupations (2.2). Game-based delivery sustains engagement, but its contribution to competence is marginal, so skill development must rest on task authenticity rather than reward structures, and randomisation must serve variety rather than reward scheduling (2.3). Career theory supports treating occupational fit as changing and as situated within a whole life, while explaining why self-report alone is a weak basis for guidance (2.4). Generative and assessment technologies make scenario production and open-response evaluation feasible, subject to explicit grounding, validated judgment, and expert oversight (2.5), and established taxonomies and benchmarks permit both grounding and evaluation without bespoke data collection (2.6). The user population imposes data protection, fairness and accuracy obligations that are design constraints rather than subsequent considerations (2.7).

Two qualifications in this evidence bear directly on what any such system may claim. The first is the marginal effect of gamification on competence\cite{r33}; the appropriate response is to locate skill development in authentic task performance and validated assessment rather than in reward mechanisms. The second is Sitzmann's finding that simulation exceeds comparison instruction only as a supplement, and underperforms as a sole method\cite{r36}; a system of this kind is therefore properly positioned as informing occupational choice and building initial self-efficacy, not as a substitute for employment or formal instruction. Both qualifications are treated in what follows as constraints on claims rather than as objections to be set aside.
